\section*{الفصل \ref{ch:g7:chemical-reactions} --- التفاعلات الكيميائية: المتفاعلات والنواتج}

\begin{solution}{exo:g7:chemical-reactions:1}
فيزيائي: انصهار الجليد، \omterm{def:g2:mixing-and-dissolving:dissolve}{وذوبان} السكر. كيميائي: \omterm{def:g4:what-a-fire-needs:burning}{احتراق} عود الثقاب،
وتحميص الخبز، وصدأ المسمار.
\end{solution}

\begin{solution}{exo:g7:chemical-reactions:2}
\omterm{def:g7:chemical-reactions:reactant}{المتفاعلات} هي الأنواع التي تُستهلك؛ \omterm{def:g7:chemical-reactions:reactant}{والنواتج} هي الأنواع الجديدة التي
تتكوّن.
\end{solution}

\begin{solution}{exo:g7:chemical-reactions:3}
كربون $+$ أكسجين $\longleftarrow$ ثنائي أكسيد الكربون.
\end{solution}

\begin{solution}{exo:g7:chemical-reactions:4}
\omterm{def:g7:chemical-reactions:reactant}{المتفاعلات}: الحديد والأكسجين. \omterm{def:g7:chemical-reactions:reactant}{الناتج}: أكسيد الحديد.
\end{solution}

\begin{solution}{exo:g7:chemical-reactions:5}
الغشاوة تبيّن الماء؛ وتعكّر ماء الجير يبيّن ثنائي أكسيد الكربون.
\end{solution}

\begin{solution}{exo:g7:chemical-reactions:6}
زنك $+$ حمض الهيدروكلوريك $\longleftarrow$ كلوريد الزنك $+$ هيدروجين.
\end{solution}

\begin{solution}{exo:g7:chemical-reactions:7}
قبل: 4 (في \omterm{def:g7:atoms-and-molecules:molecule}{جزيء} الميثان). بعد: $2 + 2 = 4$ (اثنتان في كل \omterm{def:g7:atoms-and-molecules:molecule}{جزيء}
ماء).
\end{solution}

\begin{solution}{exo:g7:chemical-reactions:8}
لا يظهر أي نوع جديد: فالملح لا يزال موجودًا، منتشرًا في الماء، ونسترجعه
\omterm{def:g3:separating-mixtures:evaporate}{بتبخير} الماء. إنه \omterm{def:g7:chemical-reactions:physical-transformation}{تحول
فيزيائي}.
\end{solution}

\begin{solution}{exo:g7:chemical-reactions:9}
معظم الخشب تفاعل مع الأكسجين فأعطى نواتج غازية، هي ثنائي أكسيد الكربون
وبخار الماء، تذهب في \omterm{def:g4:air-a-mixture-of-gases:air}{الهواء}. ولا يبقى إلا
الرماد.
\end{solution}

\begin{solution}{exo:g7:chemical-reactions:10}
هيدروجين $+$ أكسجين $\longleftarrow$ ماء. \omterm{def:g7:chemical-reactions:reactant}{والناتج} يُزرّق كبريتات النحاس
اللامائية.
\end{solution}

\begin{solution}{exo:g7:chemical-reactions:11}
قبل: $2 \times 2 = 4$ \omterm{def:g7:atoms-and-molecules:atom}{ذرات} هيدروجين وذرتا أكسجين. بعد: يحوي جزيئا
الماء $2 \times 2 = 4$ \omterm{def:g7:atoms-and-molecules:atom}{ذرات} هيدروجين و$2 \times 1 =
2$ من \omterm{def:g7:atoms-and-molecules:atom}{ذرات} الأكسجين. \omterm{def:g7:atoms-and-molecules:atom}{فالذرات} نفسها، بالأعداد نفسها.
\end{solution}

\begin{solution}{exo:g7:chemical-reactions:12}
لا يمكن \omterm{def:g7:atoms-and-molecules:atom}{لذرة} أكسجين أن تأتي من العدم: ففي الرسم ذرتا أكسجين قبل التفاعل
و3 بعده. الصحيح: \omterm{def:g7:atoms-and-molecules:atom}{ذرة} كربون واحدة $+$ \omterm{def:g7:atoms-and-molecules:molecule}{جزيء} أكسجين واحد
$\longleftarrow$ \omterm{def:g7:atoms-and-molecules:molecule}{جزيء} واحد من ثنائي أكسيد الكربون، ولا يبقى شيء زائد.
\end{solution}

\begin{solution}{pb:g7:chemical-reactions:1}
\textbf{1.} الميثان والأكسجين.

\textbf{2.} الماء.

\textbf{3.} ثنائي أكسيد الكربون.

\textbf{4.} لا: فهو لا يُستهلك، ولا يُكتب في \omterm{def:g7:chemical-reactions:word-equation}{المعادلة
اللفظية}.

\textbf{5.} ميثان $+$ أكسجين $\longleftarrow$ ثنائي أكسيد الكربون $+$
ماء.

\textbf{6.} كيميائي: يختفي نوعان (الميثان، والأكسجين)، ويظهر نوعان
جديدان تكشف عنهما الاختبارات؛ وتنطلق حرارة وضوء.

\textbf{7.} \ce{CH4}، \ce{O2}، \ce{CO2}، \ce{H2O}.

\textbf{8.} كربون واحد، و4 هيدروجين، و$2 \times 2 = 4$ أكسجين.

\textbf{9.} كربون واحد؛ و$2 \times 2 = 4$ هيدروجين؛ و$2 + 2 \times 1 = 4$
أكسجين. \omterm{def:g7:atoms-and-molecules:atom}{الذرات} نفسها بالأعداد نفسها: لم يحدث لها إلا إعادة
ترتيب.

\textbf{10.} $10 \times 2 = 20$ \omterm{def:g7:atoms-and-molecules:molecule}{جزيء} أكسجين.

\textbf{11.} 10 \omterm{def:g7:atoms-and-molecules:molecule}{جزيئات} من ثنائي أكسيد الكربون.

\textbf{12.} $10 \times 2 = 20$ \omterm{def:g7:atoms-and-molecules:molecule}{جزيء} ماء.
\end{solution}
